% Generated draft: requires mathematical and source review.
\documentclass[11pt]{article}
\usepackage[margin=25mm]{geometry}
\usepackage{fontspec}
\setmainfont{Latin Modern Roman}
\usepackage{amsmath,amssymb,mathtools,unicode-math}
\setmathfont{Latin Modern Math}
\usepackage{enumitem,needspace,xurl,hyperref}
\setlength{\parindent}{0pt}
\setlength{\parskip}{5pt}
\setlength{\emergencystretch}{3em}
\newcommand{\N}{\mathbb{N}}
\newcommand{\Z}{\mathbb{Z}}
\newcommand{\Q}{\mathbb{Q}}
\newcommand{\R}{\mathbb{R}}
\newcommand{\C}{\mathbb{C}}
\newcommand{\F}{\mathbb{F}}
\newcommand{\A}{\mathbb{A}}
\newcommand{\PP}{\mathbb{P}}
\newcommand{\E}{\mathbb{E}}
\newcommand{\eps}{\varepsilon}
\newcommand{\dd}{\,\mathrm{d}}
\begin{document}
\section*{UnsolvedMath ID: 11000152}
\textbf{Source problem code:} catalogue entry
\subsection{Definitions and mathematical statement}
A mapping-class factorization is an ordered tuple \(\mathbf t=(t_1,\ldots,t_n)\). One convention for an elementary Hurwitz move replaces adjacent factors by
\[
 (a,b)\longmapsto(b,b^{-1}ab),
\]
with its inverse also allowed; this preserves their product. Global conjugation replaces every factor by \(ht_ih^{-1}\) for one mapping class \(h\). Let \(\sim_H\) be the equivalence relation generated by these moves and simultaneous conjugation.

The question asks whether there is an algorithm that, given two factorizations in the intended mapping class group, decides whether \(\mathbf t\sim_H\mathbf t'\), terminating on both equivalent and inequivalent inputs. It also asks for useful criteria proving equivalence or inequivalence. A natural input encoding uses finite words for the factors, but the genus, marked/boundary conventions, allowed factor class and prescribed total products are not specified. Positivity, identity monodromy or a boundary multitwist are not silently imposed, nor are extra stabilization moves added.

More broadly, find useful invariants of Lefschetz fibrations derived from these factorizations, and ideally invariants of the underlying symplectic \(4\)-manifolds. The latter goal requires independence from the choice of fibration and is not automatic from Hurwitz invariance alone.

The source gives two geometric comparisons: Seidel's Fukaya-type \(A_\infty\)-category associated to a Lefschetz fibration [21], described as a promising computational approach to Lagrangian submanifolds/Fukaya categories in open exact symplectic \(4\)-manifolds; and Donaldson--Smith's enumerative invariant [9], counting embedded pseudoholomorphic curves by sections of a relative Hilbert scheme associated to a Lefschetz fibration. These examples are reported motivation, not a supplied general decision algorithm or unrestricted computability theorem.
\subsection{Short English statement}
Decide Hurwitz equivalence up to global conjugation and find effective distinguishing criteria, aiming for factorization-based Lefschetz-fibration and symplectic-manifold invariants.
\subsection{Sources}
\begin{itemize}
\item \emph{Problems on Mapping Class Groups and Related Topics}, edited by Farb, Question 2.1, supplied PDF page 141. \url{https://www.math.uchicago.edu/~farb/papers/mcgbook.pdf}. Recorded access 2026-07-29; extraction pdf-font-structure.
\item Seidel [21], Fukaya-type \(A_\infty\)-category; Donaldson and Smith [9], relative-Hilbert-scheme curve-counting invariant. Complete bibliography entries are absent. The summary also mentions Wajnryb's approach without a specific reference.
\item Local UnsolvedMath registry, \texttt{data/unsolvedmath/problems.json}, permanent ID \(11000152\), exact source code \texttt{catalogue entry}; supplied OPEN-TRIAGE summary reports general decidability unsettled.
\end{itemize}
\end{document}
